%! Author = sriadityavedantam
%! Date = 3/30/26

\section{Direct Sums and Finite Abelian Groups}

\begin{definition}[External Direct Product]
Let $H,K$ be groups.
Define
\[
G=H\times K=\{(h,k):h\in H,\ k\in K\}.
\]
Then $G$ is a group under componentwise operation.
Also,
\[
\overline{H}=\{(h,1):h\in H\}\trianglelefteq G
\]
and
\[
\overline{K}=\{(1,k):k\in K\}\trianglelefteq G.
\]
In fact,
\[
\overline{H}\cap \overline{K}=\{(1,1)\}.
\]
\end{definition}

\begin{theorem}{
Suppose $G$ is a group and $H,K\leq G$.
When is $G$ isomorphic to the direct product of $H$ and $K$.
}
This happens when $H$ and $K$ are normal in $G$, when
\[
H\cap K=\{1\},
\]
and when
\[
G=HK.
\]
\end{theorem}

\begin{remark}
Define
\[
HK\coloneqq \{hk:h\in H,\ k\in K\}.
\]
If $H,K\normal G$, then $HK$ is a subgroup of $G$.
\end{remark}

\begin{theorem}{
If $H,K\normal G$, $G=HK$, and
\[
H\cap K=\{1\},
\]
then
\[
G\cong H\times K.
\]
This is called an internal direct product.
}
Define
\[
\phi:H\times K\to G
\]
by
\[
\phi(h,k)=hk.
\]
Since $G=HK$, $\phi$ is surjective.
To show $\phi$ is a homomorphism, let $(h_1,k_1),(h_2,k_2)\in H\times K$.
Because $H$ and $K$ are normal and $H\cap K=\{1\}$, elements of $H$ commute with elements of $K$.
Indeed, for $h\in H$ and $k\in K$,
\[
hkh^{-1}k^{-1}\in H\cap K=\{1\},
\]
so $hk=kh$.
Thus
\begin{align*}
\phi((h_1,k_1)(h_2,k_2))
&=
\phi(h_1h_2,k_1k_2)\\
&=
(h_1h_2)(k_1k_2)\\
&=
h_1k_1h_2k_2\\
&=
\phi(h_1,k_1)\phi(h_2,k_2).
\end{align*}

Now compute the kernel.
If
\[
\phi(h,k)=1,
\]
then
\[
hk=1,
\]
so
\[
h=k^{-1}.
\]
Hence $h\in H\cap K$ and $k\in H\cap K$.
Since
\[
H\cap K=\{1\},
\]
we get
\[
h=k=1.
\]
Thus
\[
\ker\phi=\{(1,1)\}.
\]
So $\phi$ is injective.
Therefore $\phi$ is an isomorphism.
\end{theorem}

\begin{definition}[Disjoint Group]
If
\[
H\cap K=\{1\},
\]
we say the intersection is trivial.
\end{definition}

\begin{example}[Example 1]
Let
\[
G=\gen{x},
\qquad
\ord(x)=n.
\]
\end{example}

\begin{proof}
Suppose
\[
n=ab,
\qquad
\gcd(a,b)=1.
\]
Let
\[
H=\gen{x^a},
\qquad
K=\gen{x^b}.
\]
Then
\[
H\cong C_b,
\qquad
K\cong C_a.
\]
If $x^r\in H\cap K$, then $a\mid r$ and $b\mid r$.
Since $\gcd(a,b)=1$, this implies
\[
ab=n\mid r.
\]
Hence
\[
x^r=1.
\]
So
\[
H\cap K=\{1\}.
\]
Now let $r\in \mathbb{Z}$.
Since $\gcd(a,b)=1$, there exist $s,t\in \mathbb{Z}$ such that
\[
r=sa+tb.
\]
Then
\begin{align*}
x^r
&=
x^{sa+tb}\\
&=
x^{sa}x^{tb}\\
&=
(x^a)^s(x^b)^t.
\end{align*}
So $x^r\in HK$ for all $r$.
Hence
\[
G=HK.
\]
Therefore
\[
G\cong H\times K.
\]
\end{proof}

\begin{definition}[Direct Sum]
When $G$ is an additive abelian group, the direct product is called a direct sum.
\end{definition}

\begin{theorem*}[Fundamental Theorem]{
If $G$ is a finite abelian group, then $G$ can be written in the following equivalent ways:
\begin{enumerate}
    \item
    \[
    G\cong \mathbb{Z}_{n_1}\oplus \mathbb{Z}_{n_2}\oplus\dots\oplus \mathbb{Z}_{n_k}.
    \]
    \item
    \[
    G\cong \mathbb{Z}_{p_1^{k_1}}\oplus \mathbb{Z}_{p_1^{k_2}}\oplus \mathbb{Z}_{p_2^{\ell_1}}\oplus\dots.
    \]
    \item
    \[
    G\cong \mathbb{Z}_{d_1}\oplus \mathbb{Z}_{d_2}\oplus\dots\oplus \mathbb{Z}_{d_k},
    \]
    where the $p_i^{k_j}$ are the elementary divisors and the $d_i$ are the invariant factors.
\end{enumerate}
}
\end{theorem*}

\begin{remark}
The implication from elementary divisors to invariant factors uses the fact that if
\[
\gcd(m,n)=1,
\]
then
\[
\mathbb{Z}_{mn}\cong \mathbb{Z}_m\oplus \mathbb{Z}_n.
\]
\end{remark}

\begin{example}
Let
\[
G=\mathbb{Z}_2\oplus \mathbb{Z}_4\oplus \mathbb{Z}_4\oplus \mathbb{Z}_3\oplus \mathbb{Z}_3\oplus \mathbb{Z}_5.
\]
The primes involved are $2,3,5$.
Taking one highest available power of each prime gives
\[
2^2\cdot 3\cdot 5=60=d_3.
\]
Removing those factors leaves
\[
2^2\cdot 3\cdot 1=12=d_2.
\]
Removing again leaves
\[
2\cdot 1\cdot 1=2=d_1.
\]
So
\[
G\cong \mathbb{Z}_2\oplus \mathbb{Z}_{12}\oplus \mathbb{Z}_{60}.
\]
\end{example}

\begin{remark}
One way to begin proving the theorem is as follows.
Suppose $x_1,\dots,x_n$ generate $G$.
Define a surjective homomorphism
\[
\phi:\mathbb{Z}^n\to G
\]
by sending
\[
e_i\mapsto x_i,
\]
where
\[
e_1=(1,0,\dots,0),\quad e_2=(0,1,0,\dots,0),\quad \dots,\quad e_n=(0,0,\dots,1).
\]
Then every
\[
y=y_1e_1+\dots+y_ne_n\in \mathbb{Z}^n
\]
maps to
\[
\sum_{i=1}^n y_ix_i.
\]
If
\[
K=\ker\phi,
\]
then by the first isomorphism theorem,
\[
\mathbb{Z}^n/K\cong G.
\]
So the problem reduces to understanding subgroups of $\mathbb{Z}^n$.
\end{remark}

\begin{definition}[Free Abelian Group]
A group $G$ is a free abelian group if it has a basis.
\end{definition}

\begin{remark}
Recall from Linear Algebra.
A basis for $\mathbb{Z}^n$ is
\[
e_1,\dots,e_n,
\]
which generates $\mathbb{Z}^n$ and is linearly independent over $\mathbb{Z}$.
\end{remark}

\begin{lemma*}{
Every subgroup of $\mathbb{Z}^n$ has a basis.
}
\end{lemma*}

A direct product and direct sum are similar in spirit to bases, but now for groups.

\begin{theorem*}{
A set
\[
\{f_1,\ldots,f_m\}
\]
is a basis for a subgroup $K\leq \mathbb{Z}^n$ if every $k\in K$ can be written as
\[
k=\sum a_if_i
\]
with $a_i\in \mathbb{Z}$, and if
\[
\sum a_if_i=0
\]
implies $a_i=0$ for all $i$.
}
\end{theorem*}

We want to show that every element $k$ can be written as such a linear combination.

\begin{theorem}{
Every subgroup of $\mathbb{Z}^n$ is free abelian.
}
We prove this by induction on $n$.
Base case:
$n=1$.
Every subgroup of $\mathbb{Z}$ is cyclic.
If
\[
K\neq \{0\}
\]
is a subgroup of $\mathbb{Z}$, then
\[
K=d\mathbb{Z}
\]
for some $d\neq 0$.
So
\[
\{d\}
\]
is a basis for $K$.
Thus every subgroup of $\mathbb{Z}$ has a basis.
Now consider $n=2$.
Let
\[
K\leq \mathbb{Z}^2=\mathbb{Z}e_1+\mathbb{Z}e_2.
\]
If
\[
K\leq \mathbb{Z}e_1,
\]
then we are done by the $n=1$ case.
Assume not.
Let
\[
H\coloneqq \{b\in \mathbb{Z}:\exists y\in \mathbb{Z}\text{ such that }(y,b)=ye_1+be_2\in K\}.
\]
Then $H$ is a subgroup of $\mathbb{Z}$.
Since
\[
K\nleq \mathbb{Z}e_1,
\]
we have
\[
H\neq \{0\}.
\]
So
\[
H=d\mathbb{Z}
\]
for some $d\neq 0$.
Thus there exists $y_1\in \mathbb{Z}$ such that
\[
(y_1,d)=y_1e_1+de_2\in K.
\]
Let
\[
f_2=(y_1,d).
\]
Also,
\[
K\cap \mathbb{Z}e_1
\]
is a subgroup of $\mathbb{Z}e_1\cong \mathbb{Z}$.
If
\[
K\cap \mathbb{Z}e_1=\{(0,0)\},
\]
then one basis vector may suffice.
Otherwise,
\[
K\cap \mathbb{Z}e_1=\mathbb{Z}(a,0)
\]
for some $a\neq 0$.
Let
\[
f_1=(a,0).
\]
Then the claim is that either $\{f_2\}$ or $\{f_1,f_2\}$ is a basis for $K$.
\end{theorem}

\begin{lemma}{
If
\[
K\cap \mathbb{Z}e_1=\{(0,0)\},
\]
then $\{f_2\}$ is a basis for $K$.
If
\[
K\cap \mathbb{Z}e_1=\mathbb{Z}(a,0),
\]
then $\{f_1,f_2\}$ is a basis for $K$.
}
Let
\[
k=(z_1,z_2)\in K.
\]
Since $z_2\in H=d\mathbb{Z}$, there exists $k_2\in \mathbb{Z}$ such that
\[
z_2=dk_2.
\]
Then
\begin{align*}
k-k_2f_2
&=
(z_1,z_2)-k_2(y_1,d)\\
&=
(z_1-k_2y_1,z_2-k_2d)\\
&=
(z_1-k_2y_1,0).
\end{align*}
Thus
\[
k-k_2f_2\in K\cap \mathbb{Z}e_1.
\]
So either
\[
k=k_2f_2
\]
or
\[
k=k_1f_1+k_2f_2
\]
for some $k_1\in \mathbb{Z}$.
Hence these vectors span $K$.
Now check linear independence.
In the one-vector case, it is immediate.
In the two-vector case, let
\[
f_1=(a,0),
\qquad
f_2=(y_1,d).
\]
Suppose
\[
a_1f_1+a_2f_2=(0,0).
\]
Then
\[
(a_1a+a_2y_1,\ a_2d)=(0,0).
\]
Since $d\neq 0$, we get
\[
a_2=0.
\]
Then
\[
a_1a=0,
\]
and since $a\neq 0$, we get
\[
a_1=0.
\]
So the vectors are linearly independent.
Therefore they form a basis.
\end{lemma}